% GENERATED FILE. Edit canonical lessons, then run npm run generate:books.
% canonical-source-hash: fnv1a64:91fa470db5679cfe
% canonical-lessons: ML-C553-tottil, ML-C553-dattetukkuka, ML-C553-rantanamma, ML-C553-kutappirappu, ML-C553-sahodaran

\chapter{Orphans, Steps and Siblings}
\label{ch:ml-orphans-steps-and-siblings}
\hlchaptermodality{553}

\begin{chapteropening}
\textbf{By the end of this chapter:} \emph{I can say a cradle, to adopt, a stepmother, a sibling and a brother.}

Complete the last lesson of chapter 553: I can say a cradle, to adopt, a stepmother, a sibling and a brother.
\end{chapteropening}

\section[\texorpdfstring{\ml{തൊട്ടിൽ} (toṭṭil)}{തൊട്ടിൽ (toṭṭil)}]{\ml{തൊട്ടിൽ} (toṭṭil) — a cradle}

\label{lesson:ML-C553-tottil}



\begin{quote}
Before the new one: say the Malayalam for a boyfriend, then the Malayalam for to be in love.
\end{quote}

\subsection*{You'll want to know: \ml{തൊട്ടിൽ}}
\textbf{\ml{തൊട്ടിൽ}} — \emph{toṭṭil} — \textquotedblleft{}a cradle\textquotedblright{}.

\begin{quote}
\ml{അമ്മ} \ml{കുഞ്ഞിനെ} \ml{തൊട്ടിലിൽ} \ml{കിടത്തി}.
\end{quote}

\begin{quote}
\emph{amma kuññine toṭṭilil kiṭatti.}
\end{quote}

\begin{quote}
Mother laid the baby in the cradle.
\end{quote}

Often a cloth cradle hung from a beam; \ml{തൊട്ടിൽ} \ml{ആട്ടുക}, to rock it.

\begin{cousinweb}
\textbf{\ml{തൊട്ടിൽ}} (\emph{toṭṭil}) is native Dravidian: DEDR 3486 gives it as a cradle (DEDR 3486).
\end{cousinweb}

\begin{grammarlens}[title={What you've built}]
One more word for describing what you lived and felt.
\end{grammarlens}

\subsection*{Your turn}
\begin{itemize}
\raggedright
  \item \emph{Say it:} \emph{toṭṭil}
  \item \emph{Say it:} \emph{toṭṭil}, once more
  \item \emph{Say it:} \emph{amma kuññine toṭṭilil kiṭatti.} — Mother laid the baby in the cradle.
  \item \emph{Recall:} say the Malayalam for a boyfriend, then the Malayalam for to be in love, then say \emph{toṭṭil} again
\end{itemize}

\begin{tcolorbox}[breakable,colback=teal!4,colframe=teal!35!black,title={Before you move on}]
What does \ml{തൊട്ടിൽ} mean? (\textquotedblleft{}A cradle\textquotedblright{}.) Say it once more.
\end{tcolorbox}
\section[\texorpdfstring{\ml{ദത്തെടുക്കുക} (datteṭukkuka)}{ദത്തെടുക്കുക (datteṭukkuka)}]{\ml{ദത്തെടുക്കുക} (datteṭukkuka) — to adopt}

\label{lesson:ML-C553-dattetukkuka}



\begin{quote}
Before the new one: say the Malayalam for to be in love, then the Malayalam for a cradle.
\end{quote}

\subsection*{You'll want to know: \ml{ദത്തെടുക്കുക}}
\textbf{\ml{ദത്തെടുക്കുക}} — \emph{datteṭukkuka} — \textquotedblleft{}to adopt\textquotedblright{}.

\begin{quote}
\ml{ആ} \ml{ദമ്പതികൾ} \ml{ഒരു} \ml{കുഞ്ഞിനെ} \ml{ദത്തെടുത്തു}.
\end{quote}

\begin{quote}
\emph{ā dampatikaḷ oru kuññine datteṭuttu.}
\end{quote}

\begin{quote}
That couple adopted a baby.
\end{quote}

Past: \ml{ദത്തെടുത്തു}, \emph{datteṭuttu}.

\ml{ദത്ത്} given + \ml{എടുക്കുക} to take: to take a child as given.

\begin{cousinweb}
\textbf{\ml{ദത്തെടുക്കുക}} (\emph{datteṭukkuka}) is a Malayalam compound built from \ml{ദത്ത്}, from Sanskrit \emph{datta}, \textquotedblleft{}given\textquotedblright{} (CDIAL 6140), and \ml{എടുക്കുക}, \textquotedblleft{}to take up\textquotedblright{}, which is native (DEDR 851).
\end{cousinweb}

\begin{grammarlens}[title={What you've built}]
One more word for describing what you lived and felt.
\end{grammarlens}

\subsection*{Your turn}
\begin{itemize}
\raggedright
  \item \emph{Say it:} \emph{datteṭukkuka}
  \item \emph{Say it:} \emph{datteṭukkuka}, once more
  \item \emph{Say it:} \emph{ā dampatikaḷ oru kuññine datteṭuttu.} — That couple adopted a baby.
  \item \emph{Recall:} say the Malayalam for to be in love, then the Malayalam for a cradle, then say \emph{datteṭukkuka} again
\end{itemize}

\begin{tcolorbox}[breakable,colback=teal!4,colframe=teal!35!black,title={Before you move on}]
What does \ml{ദത്തെടുക്കുക} mean? (\textquotedblleft{}To adopt\textquotedblright{}.) Say it once more.
\end{tcolorbox}
\section[\texorpdfstring{\ml{രണ്ടാനമ്മ} (raṇṭānamma)}{രണ്ടാനമ്മ (raṇṭānamma)}]{\ml{രണ്ടാനമ്മ} (raṇṭānamma) — a stepmother}

\label{lesson:ML-C553-rantanamma}



\begin{quote}
Before the new one: say the Malayalam for a cradle, then the Malayalam for to adopt.
\end{quote}

\subsection*{You'll want to know: \ml{രണ്ടാനമ്മ}}
\textbf{\ml{രണ്ടാനമ്മ}} — \emph{raṇṭānamma} — \textquotedblleft{}a stepmother\textquotedblright{}.

\begin{quote}
\ml{രണ്ടാനമ്മ} \ml{അവനെ} \ml{സ്വന്തം} \ml{മകനെപ്പോലെ} \ml{സ്നേഹിച്ചു}.
\end{quote}

\begin{quote}
\emph{raṇṭānamma avane svantaṁ makaneppōle snēhiccu.}
\end{quote}

\begin{quote}
His stepmother loved him like her own son.
\end{quote}

\ml{രണ്ടാം} second + \ml{അമ്മ}; a stepfather is \ml{രണ്ടാനച്ഛൻ}.

\begin{cousinweb}
\textbf{\ml{രണ്ടാനമ്മ}} (\emph{raṇṭānamma}) is native Dravidian, a compound of \ml{രണ്ടാം}, \textquotedblleft{}second\textquotedblright{} (DEDR 474), and \ml{അമ്മ}, \textquotedblleft{}mother\textquotedblright{}.
\end{cousinweb}

\begin{grammarlens}[title={What you've built}]
One more word for describing what you lived and felt.
\end{grammarlens}

\subsection*{Your turn}
\begin{itemize}
\raggedright
  \item \emph{Say it:} \emph{raṇṭānamma}
  \item \emph{Say it:} \emph{raṇṭānamma}, once more
  \item \emph{Say it:} \emph{raṇṭānamma avane svantaṁ makaneppōle snēhiccu.} — His stepmother loved him like her own son.
  \item \emph{Recall:} say the Malayalam for a cradle, then the Malayalam for to adopt, then say \emph{raṇṭānamma} again
\end{itemize}

\begin{tcolorbox}[breakable,colback=teal!4,colframe=teal!35!black,title={Before you move on}]
What does \ml{രണ്ടാനമ്മ} mean? (\textquotedblleft{}A stepmother\textquotedblright{}.) Say it once more.
\end{tcolorbox}
\section[\texorpdfstring{\ml{കൂടപ്പിറപ്പ്} (kūṭappiṟappŭ)}{കൂടപ്പിറപ്പ് (kūṭappiṟappŭ)}]{\ml{കൂടപ്പിറപ്പ്} (kūṭappiṟappŭ) — a sibling}

\label{lesson:ML-C553-kutappirappu}



\begin{quote}
Before the new one: say the Malayalam for to adopt, then the Malayalam for a stepmother.
\end{quote}

\subsection*{You'll want to know: \ml{കൂടപ്പിറപ്പ്}}
\textbf{\ml{കൂടപ്പിറപ്പ്}} — \emph{kūṭappiṟappŭ} — \textquotedblleft{}a sibling\textquotedblright{}.

\begin{quote}
\ml{എനിക്ക്} \ml{മൂന്ന്} \ml{കൂടപ്പിറപ്പുകൾ} \ml{ഉണ്ട്}.
\end{quote}

\begin{quote}
\emph{enikkŭ mūnnŭ kūṭappiṟappukaḷ uṇṭŭ.}
\end{quote}

\begin{quote}
I have three siblings.
\end{quote}

Plural: \ml{കൂടപ്പിറപ്പുകൾ}, \emph{kūṭappiṟappukaḷ}.

\ml{കൂട} with + \ml{പിറപ്പ്} a birth: one born with you.

\begin{cousinweb}
\textbf{\ml{കൂടപ്പിറപ്പ്}} (\emph{kūṭappiṟappŭ}) is native Dravidian, a compound of \ml{കൂട}, \textquotedblleft{}along with\textquotedblright{} (DEDR 1882), and \ml{പിറപ്പ്}, \textquotedblleft{}a birth\textquotedblright{} (DEDR 4422).
\end{cousinweb}

\begin{grammarlens}[title={What you've built}]
One more word for describing what you lived and felt.
\end{grammarlens}

\subsection*{Your turn}
\begin{itemize}
\raggedright
  \item \emph{Say it:} \emph{kūṭappiṟappŭ}
  \item \emph{Say it:} \emph{kūṭappiṟappŭ}, once more
  \item \emph{Say it:} \emph{enikkŭ mūnnŭ kūṭappiṟappukaḷ uṇṭŭ.} — I have three siblings.
  \item \emph{Recall:} say the Malayalam for to adopt, then the Malayalam for a stepmother, then say \emph{kūṭappiṟappŭ} again
\end{itemize}

\begin{tcolorbox}[breakable,colback=teal!4,colframe=teal!35!black,title={Before you move on}]
What does \ml{കൂടപ്പിറപ്പ്} mean? (\textquotedblleft{}A sibling\textquotedblright{}.) Say it once more.
\end{tcolorbox}
\section[\texorpdfstring{\ml{സഹോദരൻ} (sahōdaran)}{സഹോദരൻ (sahōdaran)}]{\ml{സഹോദരൻ} (sahōdaran) — a brother}

\label{lesson:ML-C553-sahodaran}



\begin{quote}
Before the new one: say the Malayalam for a stepmother, then the Malayalam for a sibling.
\end{quote}

\subsection*{You'll want to know: \ml{സഹോദരൻ}}
\textbf{\ml{സഹോദരൻ}} — \emph{sahōdaran} — \textquotedblleft{}a brother\textquotedblright{}.

\begin{quote}
\ml{എന്റെ} \ml{സഹോദരൻ} \ml{വിദേശത്ത്} \ml{ജോലി} \ml{ചെയ്യുന്നു}.
\end{quote}

\begin{quote}
\emph{enṟe sahōdaran vidēśattŭ jōli ceyyunnu.}
\end{quote}

\begin{quote}
My brother works abroad.
\end{quote}

Formal or written; at home one says \ml{ചേട്ടൻ} or \ml{അനിയൻ}.

\begin{cousinweb}
\textbf{\ml{സഹോദരൻ}} (\emph{sahōdaran}) is borrowed from Sanskrit \emph{sahodara}, \textquotedblleft{}born of the same womb\textquotedblright{} (Monier-Williams), with the masculine ending -\ml{ൻ}.
\end{cousinweb}

\begin{grammarlens}[title={What you've built}]
One more word for describing what you lived and felt.
\end{grammarlens}

\subsection*{Your turn}
\begin{itemize}
\raggedright
  \item \emph{Say it:} \emph{sahōdaran}
  \item \emph{Say it:} \emph{sahōdaran}, once more
  \item \emph{Say it:} \emph{enṟe sahōdaran vidēśattŭ jōli ceyyunnu.} — My brother works abroad.
  \item \emph{Recall:} say the Malayalam for a stepmother, then the Malayalam for a sibling, then say \emph{sahōdaran} again
\end{itemize}

\begin{tcolorbox}[breakable,colback=teal!4,colframe=teal!35!black,title={Before you move on}]
What does \ml{സഹോദരൻ} mean? (\textquotedblleft{}A brother\textquotedblright{}.) Say it once more.
\end{tcolorbox}
